% Part:sets-functions-relations
% Chapter: size-of-sets
% Section: pairing-alt

\documentclass[../../../include/open-logic-section]{subfiles}

\begin{document}

\olfileid{sfr}{siz}{pai-alt}

\olsection{يوه بله جوړوونکې تابع}

\begin{explain}
د $\Nat^2$ داسې نورې شمېرنې هم شته چې د هغو د معکوس تابعو
معلومول اسانه کوي۔ يوه ئې دا ده۔ د شمېرنې پر جدولي ترتيب د
انځورولو پر ځاے، د مثبتو صحيح عددونو له داسې لړ څخه پيل وکړئ
چې په سر کښې له تشو ځايونو سره تړلے وي۔ وانګېړئ چې دا ځايونه
يو په بل پسې په لاندې ډول په جوړو $\tuple{n,m}$ ډکوو۔ لومړے
له هغو جوړو پيل وکړئ چې په لومړي مقام کښې~$0$ لري (يعنې جوړې
$\tuple{0,m}$)؛ لومړۍ جوړه (يعنې $\tuple{0,0}$) په لومړي تش
مقام کښې کېږدئ، بيا يو تش مقام پرېږدئ، دويمه جوړه (يعنې
$\tuple{0,1}$) په ورپسې تش مقام کښې کېږدئ، بيا يو مقام پرېږدئ،
او همداسې مخکښې لاړ شئ۔
\textbf{د ژباړې د نسخې سمون (PSSIZ-002):} انګرېزي نثر دويمه
جوړه په تېروتنه د درېيمې جوړې په نښه ښيي؛ لاندې اصلي جدول ئې
سمه، دويمه جوړه څرګندوي۔ زموږ د شمېرنې (نيمګړے) پيل اوس داسې دے:
\[\small
\begin{array}{@{}c c c c c c c c c c c@{}}
\mathbf 1 & \mathbf 2 & \mathbf 3 & \mathbf 4 & \mathbf 5 & \mathbf 6 & \mathbf 7 & \mathbf 8 & \mathbf 9 & \mathbf{10} & \dots \\ \\
\tuple{0,0} &  & \tuple{0,1} &  & \tuple{0,2} &  & \tuple{0,3} & & \tuple{0,4} &  & \dots \\
\end{array}
\]
همدا کار د $\tuple{1,m}$ جوړو دپاره په هغو مقامونو کښې تکرار
کړئ چې لا تش پاتې دي؛ بيا هم هر بل تش مقام پرېږدئ:
\[\small
\begin{array}{@{}c c c c c c c c c c c@{}}
\mathbf 1 & \mathbf 2 & \mathbf 3 & \mathbf 4 & \mathbf 5 & \mathbf 6 & \mathbf 7 & \mathbf 8 & \mathbf 9 & \mathbf{10} & \dots \\ \\
\tuple{0,0} & \tuple{1,0} & \tuple{0,1} &  & \tuple{0,2} & \tuple{1,1} & 
\tuple{0,3} & & \tuple{0,4} &  \tuple{1,2} & \dots \\
\end{array}
\]
په همدې طريقه $\tuple{2,m}$، $\tuple{3,m}$ او داسې نورې جوړې
ور دننه کړئ۔
\textbf{د ژباړې د نسخې سمون (OLSIZ-003):} انګرېزي نثر د دوو
پرله‌پسې قطارونو پر ځاے د يوه قطار نښه دوه ځله تکراروي؛ بشپړ
اصلي جدول د پرله‌پسې قطارونو لوستل تاييدوي۔ نو زموږ بشپړه شمېرنه
داسې پيلېږي:
\[\small
\begin{array}{@{}cc c c c c c c c c c@{}}
\mathbf 1 & \mathbf 2 & \mathbf 3 & \mathbf 4 & \mathbf 5 & \mathbf 6 & \mathbf 7 & \mathbf 8 & \mathbf 9 & \mathbf{10} & \dots \\ \\
\tuple{0,0} & \tuple{1,0} & \tuple{0,1} & \tuple{2,0}  & \tuple{0,2} & 
\tuple{1,1} & \tuple{0,3} & \tuple{3,0}  & \tuple{0,4} &  \tuple{1,2} & \dots \\
\end{array}
\]
که د پورته جدول حجرې د همدې شمېرنې له مخې وشمېرو، د زيګزاګ
پاکه کرښه به و نۀ مومو، بلکې دا ترتيب به ومومو:
\[
\begin{array}{ c | c | c | c | c | c | c | c }
& \mathbf 0 & \mathbf 1 & \mathbf 2 & \mathbf 3 & \mathbf 4 & \mathbf 5 & \dots \\
\hline
\mathbf 0 & 1 & 3 & 5 & 7 & 9 & 11 & \dots \\
\hline
\mathbf 1 & 2 & 6 & 10 & 14 & 18 & \dots & \dots \\
\hline
\mathbf 2 & 4 & 12 & 20 & 28 & \dots & \dots & \dots \\
\hline
\mathbf 3 & 8 & 24 & 40 & \dots & \dots & \dots & \dots \\
\hline
\mathbf 4 & 16 & 48 & \dots & \dots & \dots & \dots & \dots \\
\hline
\mathbf 5 & 32 & \dots & \dots & \dots & \dots & \dots & \dots \\
\hline
\vdots & \vdots & \vdots & \vdots & \vdots & \vdots & \vdots & \ddots\\
\end{array}
\]
وينو چې په~$0$ قطار کښې جوړې زموږ د شمېرنې په طاق شمېره
لرونکو مقامونو کښې دي؛ يعنې جوړه $\tuple{0,m}$ په $2m+1$ مقام
کښې ده۔ د دويم قطار جوړې، $\tuple{1,m}$، په هغو مقامونو کښې دي
چې شمېر ئې د يوه طاق عدد دوه چنده دے، په ځانګړي ډول
$2 \cdot (2m+1)$؛ د درېيم قطار جوړې، $\tuple{2,m}$، په هغو
مقامونو کښې دي چې شمېر ئې د يوه طاق عدد څلور چنده،
$4 \cdot (2m+1)$، دے؛ او همداسې نور۔ په هر قطار کښې د
$(2m+1)$ ضربي عوامل، $1$، $2$، $4$، $8$، \dots، په دقيقه توګه
د~$2$ طاقتونه دي: $1= 2^0$، $2 = 2^1$، $4 = 2^2$،
$8 = 2^3$، \dots\@ په حقيقت کښې، اړوند توان ښودونکے تل د پام
وړ جوړې لومړے غړے وي۔ نو د جوړې $\tuple{n,m}$ دپاره عامل
$2^n$ دے۔ په دې توګه عمومي فارمول لاس ته راځي:
$2^n \cdot (2m+1)$۔ خو دا له جوړو څخه \emph{مثبتو} صحيح عددونو
ته نقشه ده؛ يعنې $\tuple{0,0}$ مقام~$1$ لري۔ که وغواړو له
مقام~$0$ څخه پيل وکړو، بايد له پايلې څخه~$1$ منفي کړو۔ نو دا
لاس ته راوړو:
\end{explain}

\begin{ex}
تابع $h\colon \Nat^2 \to \Nat$ چې په لاندې ډول ورکړل شوې ده،
\[
h(n,m) = 2^n (2m+1) - 1
\]
د طبيعي عددونو د جوړو د سټ~$\Nat^2$ يوه جوړوونکې تابع ده۔
\end{ex}

\begin{explain}
له دې سره سم، زموږ د $\Nat^2$ په دويمه شمېرنه کښې جوړه
$\tuple{0,0}$ کوډ $h(0,0) = 2^0(2\cdot 0+1) - 1 = 0$ لري؛
$\tuple{1,2}$ کوډ $2^{1} \cdot (2 \cdot 2 + 1) - 1 = 2
\cdot 5 - 1 = 9$ لري؛ او $\tuple{2,6}$ کوډ $2^{2} \cdot (2
\cdot 6 + 1) - 1 = 51$ لري۔
\end{explain}

ځينې وخت د طبيعي عددونو د جوړو~$\Nat^2$ کوډول دومره بس وي،
بې له دې چې کوډونه پر هدف سټ بشپړ وي۔ د داسې کوډولو معکوسې
تابعې يوازې جزوي تابعې وي۔

\begin{ex}
تابع $j\colon \Nat^2 \to \Nat^+$ چې په لاندې ډول ورکړل شوې ده،
\[
j(n,m) = 2^n3^m
\]
يوه !!a{injective} تابع $\Nat^2 \to \Nat$ ده۔
\end{ex}

\end{document}
